\chapter{检索语义模型}

\section{认知查询表达语义意图}

面向调用方的认知查询（Cognitive Query）表达主体、领域和对象范围，包含硬约束、软偏好、语义线索、时间意图、实体与语义指向、调用方能力组合、计算预算以及所需来源信息和诊断。具体检索算法由查询计划和服务实现选择。\cite{nous-research-package}

硬约束决定候选是否合法及能否返回；软偏好只影响合法候选之间的排序。精确读取直接按身份和当前版本取值，不经过向量检索。

\section{固定查询计划与检索通道}

目标查询计划（QueryPlan）在读取候选结果前固定参与的检索通道、各通道预算、权重与融合方式、拓扑扩展预算、重排策略和材料展开预算。某通道没有产生候选时，该通道贡献为零；其余候选不会因此重新归一化。

\begin{ReportWideTable}{p{2.6cm}p{5.4cm}p{5.4cm}}
\captionof{table}{目标检索通道与验证边界}\label{tab:retrieval-lanes}
\toprule
\textbf{检索通道} & \textbf{召回职责} & \textbf{主要边界}\\
\midrule
精确读取（Exact） & 按 ID、当前版本和精确键读取。 & 不以相似度替代身份。\\
实体关系（Entity） & 认知语义指向、来源提及和关系实体。 & 分开表示三种实体语义。\\
词项检索（Lexical） & 短语、代码/路径和精确词项。 & 实际覆盖由分析器合同决定。\\
稠密检索（Dense） & 规范表示、语义片段和有来源约束的线索视图。 & 向量空间标识必须一致。\\
时间检索（Temporal） & 事件发生、主体观察、主张有效和系统记录时间。 & 不用单一创建时间代替。\\
运行时检索（Runtime） & 当前活动认知、驻留项和任务上下文。 & 与持久认知真值分开。\\
拓扑扩展（Topology） & 从可重建图进行有界扩展。 & 受主体、访问控制、生命周期和预算约束。\\
资源查询（Resource） & 查询外部知识/数据源并按需展开材料。 & 来源内容仍由外部系统负责。\\
可选重排（Rerank） & 对已固定的合法候选集进行二次排序。 & 排序结果遵循硬约束。\\
\bottomrule
\end{ReportWideTable}

\section{候选身份、多视图与跨通道融合}

候选身份由 SubjectId、领域/对象标识（owner/object\_id）和修订标识（revision\_id）构成。服务视图可再带视图标识、视图类型、序号和向量空间标识。同一修订的多个片段或线索，先在对象—修订层聚合，不因视图数量增多而获得多张独立证据票。\cite{nous-research-package}

跨通道排序可由固定查询计划给出权重：
\[
S(o)=\sum_{\ell\in L_{\mathrm{plan}}}
\frac{w_{\ell}}{k+r_{\ell}(o)}
\]
其中 $L_{\mathrm{plan}}$ 和 $w_{\ell}$ 在生成候选前确定；某对象没有某通道的排序结果时，该通道贡献为零。权重取值由后续基准评测确定。

\section{最终权威校验}

候选生成后、调用前，依据当前修订、生命周期、访问控制、抑制状态、来源有效性、时间条件和材料展开状态复核返回资格。诊断分别记录来源缺失、未召回、访问或生命周期过滤、投影过期、预算不足和外部资源不可用等原因。

\section{当前代码与目标差距}

当前代码包含精确读取、实体倒排、词项检索、稠密检索、时间、运行时、拓扑和外部资源路径，并生成检索诊断。目标设计由查询计划预先固定融合规则；当前 \texttt{\detokenize{ranking.rs}} 仍根据本批候选出现的证据类别动态计算归一化分母。研究资料包指出，该机制可能改变无关候选之间的相对顺序；固定通道集合和权重列入实现计划。\cite{nous-code-snapshot,nous-research-package}
